\id{IRSTI 52.13.25}{}

\begin{header}
\swa{S.N.~Shaposhnik, D.A.~Shokarev, N.V.~Matvienko, K.T.~Zikirova, Zh.T.~Dauletzhanova, A.I.~Konurin}{INTEGRATED GEOMECHANICAL APPROACH TO ENSURING THE SAFETY OF UNDERGROUND MINING OF ORE DEPOSITS}

\tsp{1}S.N.~Shaposhnik,
\tsp{1}D.A.~Shokarev,
\tsp{1}N.V.~Matvienko,
\tsp{1}K.T.~Zikirova,
\tsp{2}Zh.T.~Dauletzhanova\envelope,
\tsp{3}A.I.~Konurin
\end{header}

\begin{affil}
\tsp{1}LLP «Expert Pro», Ust-Kamenogorsk, Kazakhstan,

\tsp{2}K. Kulazhanov Kazakh university of technology and business, Astana, Kazakhstan,

\tsp{3}N.A. Chinakal Institute of Mining, Siberian Branch of the Russian Academy of Sciences, Novosibirsk, Russia

\corrauthor{Corresponding-author: Kaliyeva\_zhanna@mail.ru}
\end{affil}

This scientific paper considers an integrated geomechanical approach to
ensuring the stability and safety of underground mining operations,
combining an assessment of the stress--strain state of the rock mass,
rockburst (impact) hazard diagnostics, substantiation of support
parameters, the use of backfilling technologies, and the organization of
instrumental deformation monitoring. It is shown that the reliability of
design solutions is achieved through the joint use of: rock mass quality
assessments and empirical schemes for selecting support; computational
models of the stress--strain state with an adequate description of rock
mass strength; a justified determination of the initial in-situ stress
state by methods including hydraulic fracturing; and deformation
monitoring for early detection of hazardous trends. Practical
reco\-mmendations are formulated regarding the minimum dataset and the
sequence of geomechanical support throughout the life cycle of an
underground facility.

{\bfseries Keywords}: geomechanical support, excavation stability, ground
pressure management, in-situ tests, deformation monitoring.

\begin{header}
ИНТЕГРИРОВАННЫЙ ГЕОМЕХАНИЧЕСКИЙ ПОДХОД К ОБЕСПЕЧЕНИЮ БЕЗОПАСНОСТИ ПОДЗЕМНОЙ РАЗРАБОТКИ РУДНЫХ МЕСТОРОЖДЕНИЙ

\tsp{1}С.Н.~Шапошник,
\tsp{1}Д.А.~Шокарев,
\tsp{1}Н.В.~Матвиенко,
\tsp{1}К.Т.~Зикирова,
\tsp{2}Ж.Т.~Даулетжанова\envelope,
\tsp{3}А.И.~Конурин
\end{header}

\begin{affil}
\tsp{1}ТОО «Expert Pro», Усть-Каменогорск, Казахстан,

\tsp{2}Казахский университет технологии и бизнеса имени К. Кулажанова, Астана, Казахстан,

\tsp{3}Федеральное Государственное Бюджетное Учреждение Науки Институт Горного Дела им. Н.А. Чинакала СО РАН, Новосибирск, Россия,

e-mail: Kaliyeva\_zhanna@mail.ru
\end{affil}

В данной научной статье рассмотрен интегрированный геомеханический
подход к обеспечению устойчивости и безопасности подземных горных работ,
объединяющий оценку напряжённо-деформированного состояния (НДС) массива,
диагностику удароопасности, обоснование параметров крепи, применение
закладочных технологий и организацию инструментального мониторинга
деформаций. Показано, что надёжность проектных решений достигается при
совместном использовании: оценок качества массива и эмпирических схем
назначения крепи; расчётных моделей НДС с адекватным описанием прочности
массива; обоснованного определения исходного напряжённого состояния
методами (включая гидроразрыв); мониторинга деформаций для раннего
выявления опасных тенденций. Сформулированы практические рекомендации по
минимальному набору данных и последовательности геомеханического
сопровождения в жизненном цикле подземного объекта.

{\bfseries Ключевые слова}: геомеханическое сопровождение, устойчивость
выработок, управление горным давлением, натурные испытания, мониторинг
деформаций.

\begin{header}
КЕН ОРЫНДАРЫН ЖЕРАСТЫ ӘДІСПЕН ИГЕРУ ҚАУІПСІЗДІГІН ҚАМТАМАСЫЗ ЕТУГЕ АРНАЛҒАН ИНТЕГРАЦИЯЛАНҒАН ГЕОМЕХАНИКАЛЫҚ ТӘСІЛ

\tsp{1}С.Н.~Шапошник,
\tsp{1}Д.А.~Шокарев,
\tsp{1}Н.В.~Матвиенко,
\tsp{1}К.Т.~Зикирова,
\tsp{2}Ж.Т.~Даулетжанова\envelope,
\tsp{3}А.И.~Конурин
\end{header}

\begin{affil}
\tsp{1}«Expert Pro» ЖШС, Өскемен, Қазақстан,

\tsp{2}Қ.Құлажанов атындағы Қазақ технология және бизнес университеті, Астана, Қазақстан,

\tsp{3}Н.А. Ресей ғылым академиясының Сібір бөлімшесі Қытайқал тау-кен институты, Новосибирск, Ресей,

e-mail: Kaliyeva\_zhanna@mail.ru
\end{affil}

Бұл ғылыми мақалада жерасты тау-кен жұмыстарының орнықтылығы мен
қауіпсіздігін қамтамасыз етуге бағытталған, тау жыныстары массивінің
кернеулі-деформацияланған күйін (КДК) бағалауды, соққы қауіптілігін
диагностикалауды, бекітпе параметрлерін негіздеуді, закладка (толтырма)
технологияларын қолдануды және деформацияларды аспаптық мониторингтеуді
біріктіретін интеграцияланған геомеханикалық тәсіл қарастырылған.
Жобалық шешімдердің сенімділігі мына компоненттерді бірлесіп қолданғанда
қамтамасыз етілетіні көрсетілді: массив сапасын бағалау және бекітпені
тағайындаудың эмпирикалық сұлбалары; массивтің беріктігін барабар
сипаттайтын кернеулі-деформацияланған күйді есептік модельдеу; бастапқы
кернеулі күйді әдістемелік тұрғыдан негізделген анықтау (гидроажырау
әдісін қоса алғанда); қауіпті үрдістерді ерте айқындауға арналған
деформациялар мониторингі. Жерасты объектісінің өмірлік циклі бойындағы
геомеханикалық сүйемелдеудің минималды деректер жиынтығы мен кезеңдік
орындалу реттілігі бойынша практикалық ұсынымдар тұжырымдалды.

{\bfseries Түйін сөздер}: геомеханикалық сүйемелдеу, тау-кен қазбаларының
орнықтылығы, тау қысымын басқару, натуралық сынақтар, деформациялар
мониторингі.

\begin{multicols}{2}
\textbf{Introduction.} As mining advances to greater depths, the
stability of underground openings is increasingly determined by the
combined effects of fracturing, heterogeneity, scale effects, and high
levels of natural stresses. Under these conditions, a design solution
based on a single method becomes vulnerable: an engineering rock mass
classification provides a rapid assessment but does not describe stress
redistribution; numerical modeling enables analysis of the
stress--strain state but depends on the correctness of the input
parameters; in-situ measurements and monitoring provide objective
feedback but require interpretation within a unified framework. In
practice, a reliable methodology is one in which rock mass quality
assessment, determination of the stress state, computational analyses,
and deformation control are integrated into a closed-loop support cycle
for mining operations.

The aim of this work is to substantiate an integrated approach to ground
pressure management and the stability of underground excavations, and to
formulate a sequence of actions that ensures consistency between
engineering assessments of the rock mass, computational models, and
in-situ control data {[}1--3{]}.

The scientific problem addressed in this study consists in the lack of a
unified methodological framework that integrates empirical rock mass
classification, numerical modeling, and in-situ monitoring into a
consistent decision-making system for underground mining conditions.

The research hypothesis assumes that the reliability of geomechanical
support can be significantly improved by implementing a closed-loop
integrated approach, where empirical classification, stress--strain
modeling, and monitoring data are consistently combined and iteratively
refined.

The objectives of the study are:

- to analyze the limitations of existing geomechanical approaches;

- to develop a structured integrated methodology combining empirical,
analytical and

observational methods;

- to substantiate the sequence of geomechanical support stages;

to validate the applicability of the proposed approach using comparative
analysis of field and calculated data.

\emph{Scientific novelty}

The scientific novelty of the study lies in:

- the formalization of an integrated geomechanical framework combining
empirical, analytical and observational approaches into a unified
system;

- the development of a closed-loop methodology for ground pressure
management based on continuous feedback between modeling and monitoring;

- the substantiation of a sequential algorithm for geomechanical support
throughout the life cycle of underground excavations;

- the integration of rockburst hazard criteria (K1, K2) into a unified
interpretation scheme for engineering decision-making.

\textbf{Materials and methods}. The research methodology is based on the
integration of three main components:

- еmpirical assessment of rock mass quality using RQD, RMR, Q-system and
GSI for initial classification and zoning;

- аnalytical and numerical modeling of the stress--strain state using
rock mass strength criteria (Hoek--Brown) and deformation parameters
derived from empirical correlations;

- insitu measurements and monitoring including hydraulic fracturing
for stress determination and instrumental monitoring (e.g.,
extensometry, TDR).

The methodological framework follows a closed-loop approach:

Rock mass classification → parameter estimation → numerical modeling →
in-situ validation → model correction.

This ensures consistency between predicted and observed behavior.

Modern underground geomechanics evolves at the intersection of three
lines: engineering classification of the rock mass, computational
modeling, and observational control. Historically, empirical methods
gained broad acceptance first, because they provide a formalized
language for describing the rock mass and make it possible to make
decisions under inevitably incomplete initial information. One of the
key elements of field description became the rock quality designation
index, which translates ``visual'' fracturing into a quantitative
indicator and thereby improves comparability of assessments between
sites and observers {[}3{]}.

Subsequently, empirical methods developed toward rating systems that
sought to account not only for the frequency of discontinuities, but
also for their engineering significance: surface condition, infilling,
weathering, water influence, and the orientation of discontinuities
relative to the excavation. This approach enabled a transition from
simple ``description'' to engineering interpretation, where the final
rating represents an approximate measure of rock mass quality in the
context of excavation stability and support selection {[}2,4{]}. A
significant step was the development of systems oriented toward support
design. In works on rock mass classification for tunnel support design
it was shown that rock mass quality can be described as a function of
several parameters reflecting blockiness, discontinuity properties, and
the influence of water, and then the resulting index can be used to
select the support type and intensity {[}1{]}.

This logic was further developed in directions associated with the
Norwegian design school, where the classification index is used as a
basis for choosing a combination of rock bolts and shotcrete, with the
possibility of updating the assessment as the rock mass is exposed
{[}5{]}. However, empirical methods have fundamental limitations. First,
the final rating depends on the quality of field description and on the
consistency of criteria; therefore, methodological discipline and
personnel training are required. Second, rating systems inevitably
average the rock mass and poorly reflect local ``weak spots'': single
faults, zones of tectonic reworking, or abrupt changes in lithology.
Third, as depth and stress level increase, stability becomes determined
not so much by block geometry as by stress redistribution and the
development of shear and tensile zones around openings. Therefore, as a
rule, an empirical assessment should be treated as the first level of
analysis that defines the initial classification picture and helps rank
sites by expected complexity, but does not replace computation and
in-situ verification {[}2{]}.

In literature it is emphasized that empirical methods should be used not
as a ``tabular replacement of calculations'', but as a tool for primary
diagnosis and comparison of sites. In particular, an overview of
empirical methods for assessing rock mass condition notes the need to
correctly choose an approach for specific conditions and to recognize
that any rating system is ultimately a model with a limited domain of
applicability {[}2{]}. This position is consistent with international
practice, where empiricism serves as an ``input'' for computational
parameterization and for observational control.

In underground excavation design, a key problem is the transition from
laboratory characteristics of intact rock to rock mass parameters. In
the laboratory, the strength and deformability of a monolithic specimen
are measured, whereas a rock mass is a system of blocks and
discontinuities, and its behavior is determined by the joint action of
the matrix and weakness surfaces. To describe rock mass strength, an
empirical failure criterion is widely used that accounts for the
transition from intact rock to a rock mass by means of the Geological
Strength Index. In the early-2000s edition, this criterion became
widespread as a unified framework enabling estimation of rock mass
strength for different degrees of disturbance and the use of parameters
in engineering calculations {[}6{]}.

This logic was further developed in directions associated with the
Norwegian design school, where the classification index is used as the
basis for selecting a combination of rock bolting and shotcrete, with
the possibility of updating the assessment as the rock mass is exposed
{[}6{]}. A separate development line is represented by empirical methods
specially adapted to mining conditions, where technological factors are
present: the influence of blasting, stress changes during extraction,
and effects of discontinuity orientation relative to stopes and pillars.
For such tasks, a mining rating system was proposed in which the basic
rock mass rating is corrected for mining factors; this makes the
empirical description more ``applied'' for design at mining operations
{[}7{]}.

The Geological Strength Index is valuable in that it formally links
geological description with computational parameterization. It is
``geologically friendly'' because it is oriented toward the observable
structure and the condition of discontinuities rather than only
numerical indices. In the classical work devoted to this index, it is
shown that correct assessment of the structural type and the condition
of discontinuity surfaces is decisive for deriving rock mass design
parameters {[}8{]}. At the same time, it is important to understand that
parameterization based on the Geological Strength Index and the rock
mass failure criterion does not eliminate the need for calibration using
in-situ manifestations. Even with careful structural assessment, there
remain factors that are difficult to capture ``from a picture'': the
degree of unloading and damage due to blasting, residual stresses, and
local tectonic zones. Therefore, in applied geomechanics, an approach is
common in which rock mass damage is introduced through a separate
coefficient reflecting technological influence and the quality of mining
practice. This idea is methodically formalized in studies devoted to
empirical estimation of the rock mass deformation modulus, where along
with the Geological Strength Index a disturbance factor is used to
account for technology-related effects {[}9{]}.

To quantitatively describe the deformation behavior of rocks and to
subsequently assess brittleness, a ``stress--strain'' diagram obtained
under controlled conditions of uniaxial or triaxial compression is used.
Below, the test mode provisions and the processing procedure for the
experimental diagram are presented.

In Fig. 1, the deformation curve is shown. The test mode consists in
maintaining a constant longitudinal strain rate of the rock specimen
during loading, which is achieved when using a servo-controlled press
operating in an automated mode of controlled longitudinal strains. Tests
are carried out in uniaxial or triaxial compression with a strain-rate
range from 10\textsuperscript{-3} to 10\textsuperscript{-8}
s\textsuperscript{-1}. From the specimen deformation curve (Fig. 1), the
elastic modulus E and the post-peak modulus M, the total strain
ε\textsubscript{total}, and the elastic strain ε\textsubscript{e} are
determined.

\fig[0.8\columnwidth]{g/image109}[Fig.1 - Schematic ``stress--strain'' diagram under
uniaxial compression: σ\tsb{1}---axial stress on the specimen, MPa;
ε\tsb{1} - longitudinal strain of the specimen; E - elastic modulus of
the specimen, MPa; M-post-peak modulus of the specimen, MPa; σ\tsb{re
s}- residual strength, MPa; σ\tsb{uc} - uniaxial compressive strength,
MPa]

The figure is provided for the visual interpretation of the relationship
discussed in the text.

To move from analysis of the diagram shape to quantitative
classification, indicators are used that relate elastic properties to
the character of post-peak softening. As calculation criteria, the
ratio-of-moduli coefficient and the elastic-strain fraction coefficient
are applied.

When assessing brittleness (rockburst hazard) of rocks, the following
coefficients are considered:

\begin{equation*}
K_{1} = \frac{E}{M}
\end{equation*}

\begin{equation*}
K_{2} = \frac{\varepsilon_{e}}{\varepsilon_{\text{total}}}
\end{equation*}

For K\tsb{1} <{} 1 the rock is considered
rockburst-prone; for K\tsb{1} \textgreater{} 1 it is
considered non-rockburst-prone. For K\tsb{2} \textgreater{}
0.7 the rock is rockburst-prone; for K\tsb{2} <{} 0.7
the rock is considered non-rockburst-prone. If at least one coefficient
confirms rockburst proneness, a conclusion is made that the rock is
rockburst-prone.

Thus, criteria K\tsb{1} and K\tsb{2} provide a
unified approach to interpreting test results and make it possible to
compare rockburst proneness under different loading conditions. Below,
the requirements for sample size and data processing are presented.

Rock mass deformability is one of the most sensitive parameters in
stability problems. The rock mass deformation modulus largely determines
the expected convergence of excavations, stress redistribution, and the
tendency for the development of damaged zones around an opening. Direct
estimation of the rock mass deformation modulus from in-situ tests is
often constrained by cost and organizational conditions; therefore,
empirical relationships are widely used in practice. A systematic
discussion of such relationships and a formula proposed in the mid-2000s
that accounts for the Geological Strength Index, intact rock modulus,
and a disturbance factor have become one of the most cited bases for
engineering estimates {[}9{]}. However, comparative studies show that
different empirical methods can yield substantially different values of
the rock mass deformation modulus even for identical initial rock mass
characteristics. In a study based on dozens of rock mass scenarios, it
is emphasized that the choice of correlation should be determined by the
rock mass type and the expected deformation mechanism; otherwise,
computational predictions may systematically overestimate or
underestimate deformations {[}10{]}. This conclusion has a direct
practical implication: it is advisable to estimate the rock mass
deformation modulus by several methods and then verify consistency with
in-situ observations and instrumental monitoring.

Determination of deformation properties of rocks is carried out in
accordance with GOST 28985--91. The standard applies to hard rocks with
uniaxial compressive strength not less than 5 MPa and establishes the
method for determining their deformation characteristics. Longitudinal
ε\tsb{1} and transverse ε\tsb{2} strains, axial load
(P), and lateral pressure are recorded automatically by the measuring
system and saved to an Excel file. The loading rate is 0.1 mm/min until
specimen failure. Based on the measurements, ``stress σ -- strain ε''
relationships are constructed (Fig.2). Stress is calculated from the

\begin{equation*}
\sigma = \frac{P}{A}
\end{equation*}

\begin{equation*}
\sigma = \frac{4P}{\pi d^{2}}
\end{equation*}

The deformation modulus Ed and the Poisson's ratio ν are determined in
the selected stress range (σ\_start -- σ\_end) of the linear portion of
the loading branch using the formulas:
expression:

\begin{equation*}
E_{d} = \frac{\sigma_{\text{end}} - \sigma_{\text{start}}}{\varepsilon_{1\text{end}} - \varepsilon_{1\text{start}}}
\end{equation*}

\begin{equation*}
\nu = \frac{\epsilon_{2\text{end}} - \epsilon_{2\text{start}}}{\epsilon_{1\text{end}} - \epsilon_{1\text{start}}}
\end{equation*}

where σ\tsb{end}, σ\tsb{start} are the stress values
at the end and at the beginning of the loading range;
ε\tsb{1end}, ε\tsb{1start} are the relative
longitudinal strains at the end and at the beginning of the range;
ε\tsb{2end}, ε\tsb{2start} are the relative
transverse strains at the end and at the beginning of the range. In the
case of a nonlinear ``stress--strain'' relationship, a linear
approximation method is applied.
\end{multicols}

\fig[0.5\textwidth]{g/image111}[Fig.2 - ``stress σ -- strain ε'' relationships for determining
deformation properties of rocks]

\begin{multicols}{2}
As depth and geological complexity increase, the role of natural
stresses becomes decisive. Under such conditions, an engineering rock
mass classification focused on fracturing and discontinuity quality does
not provide sufficient information for predicting stress-controlled
manifestations: slabbing, spalling, bursting, and dynamic failure.
Therefore, in-situ methods for stress determination are applied in
stressed areas. One of the most common is hydraulic fracturing, as well
as hydraulic testing of pre-existing fractures. The methodological bases
and limitations of these methods are described in detail in the
recommended methods of the International Society for Rock Mechanics. It
is emphasized that results are sensitive to borehole orientation, the
presence and nature of natural fractures, borehole wall condition, and
the correctness of interpreting breakdown and shut-in pressures. It is
especially important that a reliable assessment requires a series of
tests and comparison of results, because the rock mass is generally
heterogeneous and single tests may reflect a local feature rather than
the regional stress field {[}11{]}. The stress--strain relationships
presented in Figures 1 and 2 correspond to generalized deformation
behavior of rocks, including crack closure, elastic deformation, stable
crack propagation and post-peak failure stages, which are widely
confirmed in recent studies {[}12,13{]}. Modern experimental and
numerical studies also demonstrate that rock behavior under loading is
nonlinear and includes distinct deformation stages reflected in
stress--strain curves {[}14,15{]}.

A separate layer of literature is devoted to the fact that the stress
state and deformation mechanism of rock masses under complex
natural--technogenic processes depend on the combined influence of
geology and technogenic impact. Using karst sinkholes as an example, it
is shown that deformation development is staged and requires an
integrated approach in which in-situ observations are linked to a model
representation of stress redistribution {[}16{]}. This result is
important for underground mines as well: it confirms that stress-related
manifestations rarely have a single cause, and correct interpretation
requires systematic comparison of data.

For underground ore mining, methods oriented toward stope and excavation
stability play an important role. A special place is occupied by the
stability graph method, which relates a ``modified stability number'' of
the rock mass to excavation geometry and allows prediction of the
likelihood of collapse or the need for enhanced support. Contemporary
publications emphasize that this method is widely used in underground
polymetallic mines and is continuously refined as practical data
accumulate {[}17{]}. Crucially, such methods are observational by
nature: they work well when an operation can collect and classify actual
stability data and correctly relate them to rock mass parameters.
Therefore, in engineering practice it is rational to apply stability
graphs in combination with rock mass classification, stress--strain
modeling, and deformation monitoring. In that case, the empirical graph
is not a ``replacement'' but part of a comprehensive toolkit.

Support of underground openings in mines has evolved from simple ``load
bearing'' solutions to systems where not only static strength matters,
but also deformation compatibility with the rock mass, the ability to
hold blocks, accommodate movements, and under certain conditions resist
dynamic impacts. Classical works on rock mechanics for underground
operations emphasize that the support design must correspond to the
failure mechanism and be part of the overall stability design concept
{[}18{]}. A significant contribution to support practice was made by
manuals on cable bolting, which describe design principles, installation
technology, quality control, and observational adjustment based on
instrumentation data. These works emphasize that support should be
designed iteratively: the initial design is refined based on actual rock
mass behavior and the characteristics of installed support {[}19{]}.

For conditions prone to dynamic manifestations, the literature
separately considers principles of designing support capable of
withstanding dynamic loads and large displacements. It is shown that
support selection under such conditions should be iterative and must
include verification based on observational data, since static
approaches do not reflect the real system behavior under dynamic failure
{[}20{]}. In domestic studies devoted to assessing rock mass condition
at large mines, it is also emphasized that classification assessments
and observations should be combined with consideration of real
manifestations and technological factors {[}11{]}. In works on
integrated high-productivity extraction technologies, emphasis is placed
on the fact that stability and safety must be considered as part of the
technological mining scheme, and geomechanical support becomes a key
factor in ensuring productivity {[}21{]}.

The observational concept in geomechanics is that the design decision is
treated as a hypothesis that is refined as data on rock mass behavior
are obtained. Therefore, monitoring should not be limited to recording
``consequences''; its task is early detection of deformation trends and
localization of active zones. In this context, instrumental methods are
important that make it possible to obtain information on the depth and
position of shear zones. For underground conditions, it is essential
that instrumental data become meaningful only when interpreted jointly
with a stress--strain model and with an understanding of rock mass
structure. In other words, monitoring is the ``language of feedback'',
but the decision is formed based on reconciliation of: classification
description, stress model, computed rock mass parameters, and actual
measurements.

Time domain reflectometry is a method in which the reflection of an
electromagnetic pulse in a cable line is used to determine the location
of damage or local change corresponding to shear or deformation in the
rock mass. The methodology and examples of its use to determine the
depth of a shear zone are described in detail in the classic work on
slope monitoring; the key part of the approach transfers well to the
tasks of localizing deformations in geotechnical systems provided
correct installation and interpretation of signals {[}22{]}.

Summarizing the literature, a stable pattern can be identified: each
individual approach solves only part of the problem. Empirical methods
rank the rock mass well by quality and help select support type, but
they reflect stress-controlled failure mechanisms less effectively. Rock
mass parameterization through the Geological Strength Index and a rock
mass failure criterion creates a bridge to computations but requires
correct structural assessment and calibration. In-situ stress
determination methods provide fundamentally important information but
are sensitive to heterogeneity and require a series of tests. Monitoring
and support testing provide feedback but need interpretation. Therefore,
an integrated approach in which these elements are combined into a
closed support procedure represents the most substantiated methodology
for underground ore mines, especially under high stresses, variable rock
mass structure, and increased productivity requirements {[}23-25{]}.

{\bfseries Results and discussion.} Generalization of in-situ observations
and computational interpretations shows that the stability of
underground excavations is most correctly described by a consistent set
of characteristics including rock mass quality, stress state parameters,
and confirmation of support effectiveness by in-situ data.

Empirical methods for assessing rock mass condition provide a
reproducible engineering language for comparing sites and make it
possible to identify zones where stability reduction is expected when
excavation geometry or mining technology changes. Their strength lies in
speed and applicability under limited initial data. However, for
stressed and structurally heterogeneous rock masses, empirical indices
should be treated as a preliminary assessment that requires confirmation
by computational analysis and in-situ stress measurements.

In-situ determination of the stress state is especially important for
areas where instability manifestations are stress-controlled or
associated with stress redistribution due to mining. Studies of the
stress state and deformation mechanisms of rock masses under complex
conditions, including natural--technogenic sinkhole zones, show that
deformation processes depend on the combination of geological
heterogeneity and technogenic factors, and interpretation requires
linking field data and modeling. In this context, in-situ stress
determination methods should be applied in series, and design solutions
should be checked against the range of admissible initial parameters,
which is consistent with recommended methods and reduces the risk of
errors caused by local rock mass variability.

Quantitative results and relationships obtained from processing in-situ
data and empirical parameterization are presented below. An example of
quantitative assessment of rock mass quality and the associated expected
deformability for two ore bodies is then provided.

Parameterization of rock mass strength and deformability for
computational analyses should be based on geologically justified indices
and should account for scale effects. The use of a rock mass failure
criterion and empirical estimates of the deformation modulus makes it
possible to form design parameters consistent with the observed
structural state. Comparative studies show that discrepancies between
empirical deformability estimates can be significant; therefore, the
correctness of the correlation choice should be confirmed by comparison
with in-situ deformation observations and excavation behavior.

In-situ verification of support and deformation monitoring closes the
support cycle. Support testing makes it possible to assess actual
load-bearing capacity and technological reproducibility, while
monitoring records the transition from quasi-static behavior to the
development of local deformation zones. Practice in applying domain
reflectometry shows its effectiveness for identifying the position and
depth of active shear zones, which increases the information content of
geomechanical control and allows decisions to be made before hazardous
deformations develop. As a result, the integrated approach enables a
transition from reacting to consequences to risk management based on
data and their consistent interpretation, which is particularly
important for high-productivity underground mining technologies.
\end{multicols}

\tcap{Table 1 - Rock mass quality indicators and deformability assessment}
\begin{longtblr}[
  label = none,
  entry = none,
]{
  width = \linewidth,
  colspec = {Q[100]Q[63]Q[183]Q[163]Q[135]Q[117]Q[188]},
  cells = {c, font = \small},
  row{1} = {font=\bfseries\small},
  hlines,
  vlines,
}
Site & Ore body & Rock Quality Designation (RQD), \% & Geological Strength Index (GSI) & Rock Mass Rating (RMR) & Q-system index (Q) & Rock mass deformation modulus, GPa\\
Quartzite Hills & 1 & 55.8 & 50 & 54 & 4.5 & 16.51\\
Quartzite Hills & 4 & 35.1 & 44 & 49 & 3.3 & 14.89
\end{longtblr}

To interpret rock mass behavior, it is important to consider mechanical
properties of the main lithological varieties; summary values are
provided below.

\tcap{Table 2 - Averaged mechanical properties of the main lithological varieties}
\begin{longtblr}[
  label = none,
  entry = none,
]{
  width = \linewidth,
  colspec = {Q[150]Q[110]Q[163]Q[155]Q[171]Q[171]Q[94]},
  cells = {c, font = \small},
  row{1} = {font=\bfseries\small},
  hlines,
  vlines,
}
Lithology & Tensile strength, MPa & Uniaxial compressive strength, MPa & Hoek–Brown parameter mi & Intact rock deformation modulus, GPa & Young’s modulus of specimen, GPa & Poisson’s ratio\\
Tuffs & 6.7 & 76.5 & 11 & 11.45 & 51.14 & 0.25\\
Sedimentary rocks & 8.1 & 48.1 & 7 & 10.95 & 40.13 & 0.25\\
Metasomatites & 11.1 & 45.6 & 7 & 9.00 & 89.2 & 0.25
\end{longtblr}

Below are the results of a series of in-situ tests illustrating the
variability of the first loading pressure under similar conditions.
\vspace{-1em}
\tcap{Table 3 - In-situ stress testing: first loading pressure and distance}
\begin{longtblr}[
  label = none,
  entry = none,
]{
  cells = {c, font = \small},
  row{1} = {font=\bfseries\small},
  hlines,
  vlines,
}
Interval & Distance, m & First loading pressure, MPa\\
1 & 10 & 29.9\\
1 & 9 & 23.3\\
2 & 10 & 24.2\\
2 & 9 & 23.7\\
3 & 8 & 29.1\\
3 & 7 & 16.5
\end{longtblr}

The plot makes it possible to visually assess the scatter of first
loading pressures and its relationship with the distance to the opening
zone in the test series (Figure 3).

\begin{figs}
\fig[0.45\textwidth]{g/image112}[Fig.3 - First loading pressure and distance to the opening zone]
\fig[0.45\textwidth]{g/image113}[Fig.4 - Relationship between pressure and load in
rock-bolt testing]
\end{figs}

To verify support reliability, the results of roof rock-bolt testing
from several measurements are provided.

\tcap{Table 4 - Rock-bolt tests}
\begin{longtblr}[
  label = none,
  entry = none,
]{
  cells = {c, font = \small},
  row{1} = {font=\bfseries\small},
  hlines,
  vlines,
}
Location & Pressure, MPa & Load, t\\
Roof rock & 42 & 20.2\\
Roof rock & 45 & 21.8\\
Roof rock & 43 & 20.8\\
Roof rock & 44 & 21.2
\end{longtblr}

\begin{multicols}{2}
The plot shows the relationship between injection pressure and measured
load, which is used to check reproducibility of the test results.

Figure 4 is provided as a visual representation of the relationship used
in interpretation and discussion of results.

The results obtained demonstrate that:

- discrepancies between empirical deformation modulus estimation methods
can reach significant values, confirming the need for multi-method
validation;

- integrated interpretation of in-situ stress measurements and
monitoring data reduces uncertainty in geomechanical predictions;

- the application of the proposed integrated framework improves the
reliability of support design decisions under complex geomechanical
conditions.

{\bfseries Conclusions}. The integrated approach that combines empirical
rock mass assessment, selection of strength and deformation parameters
for calculations, and in-situ stress determination and monitoring-based
verification of decisions is a justified basis for geomechanical support
of underground operations. Empirical methods are useful for ranking
sites and for preliminary selection of excavation support. However,
under high stress levels and rock mass heterogeneity, they should be
supplemented by in-situ stress measurements and computational
verification of the stability of the adopted decisions. Design
parameters of the rock mass must be checked against actual observations:
deformations and the real behavior of excavations. This is important
because different empirical relationships can produce noticeably
different deformability estimates for the same initial information.
Deformation monitoring and in-situ verification of support performance
should be a mandatory part of ground pressure management. They provide
feedback and make it possible to adjust design and technological
decisions in time. The proposed integrated approach can be considered as
a methodological basis for modern geomechanical support systems,
particularly under conditions of increasing mining depth and stress
levels.
\end{multicols}

\begin{center}
{\bfseries Литература}
\end{center}

\begin{refs}
1. Макаров А.Б. Управление горным давлением при подземной
разработке месторождений. Системы разработки/SRK, Москва. -2021. -73
с.

2. Абрамкин Н.И., Ефимов В.И., Мансуров П.А. Эмпирические методики оценки
состояния массива горных пород// Известия УГГУ. -2021. -T.4(64). -С.
109--115. DOI 10.21440/2307-2091-2021-4-109-115.

3. Deere D.U., Deere D.W. The Rock Quality Designation (RQD) Index in
Practice/Rock Classification Systems for Engineering Purposes. ASTM
STP 984. -1988. ISBN 0-8031-0988-1.

4. Bieniawski Z.T. Engineering Rock Mass Classifications: A Complete
Manual for Engineers and Geologists in Mining, Civil, and Petroleum
Engineering// New York: Wiley. -1989. ISBN:~9780471601722, 0471601721.

5. Grimstad E., Barton N. Updating of the Q-system for the Norwegian
Method of Tunnelling// Engineering Classification of Rock Masses for
the Design of Tunnel Support. Rock Mech. -1993. -P.189-239.

6. Hoek E., Carranza-Torres C., Corkum B., et al. Hoek--Brown Failure
Criterion - 2002 Edition // Proceedings of NARMS-TAC conference.
-2002. -P.267--273.

7. Laubscher D.H. A Geomechanics Classification System for the Rating of
Rock Mass in Mine Design// Journal of the South African Institute of
Mining and Metallurgy. -1990. -Vol.90(10). -P.257-273. DOI
10.10520/AJA0038223X\_1954.

8. Marinos P., Hoek E. Geological Strength Index: A Geologically Friendly
Tool for Rock Mass Strength Estimation// Proceedings of GeoEng 2000.
-2000. -P.1422-1446.

9. Hoek E., Diederichs M.S. Empirical Estimation of Rock Mass Modulus //
International Journal of Rock Mechanics and Mining Sciences. -2006.
-Vol.43(2). -P.203-215. DOI 10.1016/j.ijrmms.2005.06.005.

10. Polemis Júnior K., Chagas da Silva Filho F., Pinheiro Lima-Filho F.
Estimating the Rock Mass Deformation Modulus: A Comparative Study of
Empirical Methods Based on 48 Rock Mass Scenarios// REM - International
Engineering Journal. -2021. -Vol.74(1). -P.39-49. DOI
10.1590/0370-44672019740150.

11. Еременко В.А., Айнбиндер И.И., Пацкевич П.Г., Бабкин Е.А. Оценка
состояния массива горных пород на рудниках ЗФ ОАО "ГМК "Норильский
никель" // Горный информационно-аналитический бюллетень. -2017. -№1.
-C.5-17. ISSN 0236--1493.

12. Feng X.T., Gong B., Liang Z.,~\emph{et al.} tudy of the Dynamic
Failure Characteristics of Anisotropic Shales Under Impact Brazilian
Splitting // Rock Mechanics and Rock Engineering. -2024. -Vol.57. -P.
2213-2230. DOI 10.1007/s00603-023-03673-w.

13. Wen S., et al. Mechanical behavior of layered rocks under dynamic
loading // International Journal of Rock Mechanics and Mining Sciences.
-2023. -Vol.170:105533. DOI 10.1016/j.ijrmms.2023.105533.

14. Qiang Y., et al. Study on closing and cracking stress calculation
method of fractured rock //Frontiers in Earth Science. -2023. -Vol.
10:839304. DOI 10.3389/feart.2022.839304.

15. Huang L., et al. Mathematical problems in rock mechanics and
stress--strain evolution under thermal effects // Mathematics. -2023.
-Vol.11(1):67. DOI 10.3390/math11010067.

16. \href{https://journals.eco-vector.com/0869-7809/search/authors/view?firstName=\%D0\%AE.&middleName=\%D0\%90.&lastName=\%D0\%9C\%D0\%B0\%D0\%BC\%D0\%B0\%D0\%B5\%D0\%B2}{Мамаев
Ю.А.},~\href{https://journals.eco-vector.com/0869-7809/search/authors/view?firstName=\%D0\%90.&middleName=\%D0\%9D.&lastName=\%D0\%92\%D0\%BB\%D0\%B0\%D1\%81\%D0\%BE\%D0\%B2}{Власов
А.Н.},~\href{https://journals.eco-vector.com/0869-7809/search/authors/view?firstName=\%D0\%9C.&middleName=\%D0\%93.&lastName=\%D0\%9C\%D0\%BD\%D1\%83\%D1\%88\%D0\%BA\%D0\%B8\%D0\%BD}{Мнушкин
М.Г.},~\href{https://journals.eco-vector.com/0869-7809/search/authors/view?firstName=\%D0\%90.&middleName=\%D0\%90.&lastName=\%D0\%AF\%D1\%81\%D1\%82\%D1\%80\%D0\%B5\%D0\%B1\%D0\%BE\%D0\%B2}{Ястребов
А.А.} Изучение напряженного состояния и механизма деформирования
массивов горных пород при образовании природно-техногенных карстовых
провалов// Геоэкология. Инженерная геология. Гидрогеология.
Геокриология. -2019. -Т.1. -P.46-59. DOI
\href{https://doi.org/10.31857/S0869-78092019146-59}{10.31857/S0869-78092019146-59}.

17. Potvin Y. Empirical Open Stope Design in Canada//The University of
British Columbia. - 1988. DOI 10.14288/1.0081130.

18. Brady B.H.G., Brown E.T. Rock Mechanics for Underground Mining.3rd
ed./ Dordrecht: Springer. -2006. DOI 10.1007/978-1-4020-2116-9.

19. Hutchinson D.J., Diederichs M.S. Cablebolting in Underground Mines/
BiTech Publishers. -1996. -406 p. ISBN 9780921095378.

20. Kaiser P.K., Cai M. Design of Rock Support System under Rockburst
Condition// Journal of Rock Mechanics and Geotechnical Engineering.
-2012. -4(3). -P.215-227. DOI 10.3724/SP.J.1235.2012.00215.

21. Абрамкин Н.И., Дородний А.В., Бухарбаев И.У.Анализ интегрированной
технологии высокопроизводительной отработки запасов выемочных участков
угольных шахт//Уголь. -2019. -№1(1114). -С.40-45. DOI
\href{http://dx.doi.org/10.18796/0041-5790-2019-1-40-45}{10.18796/0041-5790-2019-1-40-45}.

22. Kane W.F. Monitoring Slope Movement with Time Domain Reflectometry//
Engineering, Environmental Science, Geology. -2000.

23. Sarr D., Sall O., Ndiaye M., Lompo N. Comparative Analysis of the
Hard Hillsides Stability by Empirical Methods and Limit Equilibrium:
Case of Ultra Basic and Andesites of Mako and Marbles of Bandafassi
(Senegal)// Geomaterials. -2019. -Vol.9(3). -P.67-79. DOI
10.4236/gm.2019.93006.

24. Haimson B.C., Cornet F.H. ISRM Suggested Methods for Rock Stress
Estimation - Part 3: Hydraulic Fracturing (HF) and/or Hydraulic Testing
of Pre-existing Fractures (HTPF)// International Journal of Rock
Mechanics and Mining Sciences. -2003. -Vol.40(7-8). -P.1011-1020. DOI
10.1016/j.ijrmms.2003.08.002.

25. Kuranov A.D., Bagautdinov I.I. Experience of Geomechanical Research
and Calculations toward Ore Deposit Development in Complex Mining
Conditions// ~ISRM European Rock Mechanics Symposium - EUROCK 2018.
-2018.
\end{refs}

\begin{center}
{\bfseries References}
\end{center}

\begin{refs}
1. Makarov A.B. Upravlenie gornym davleniem pri podzemnoj razrabotke
mestorozhdenij. Sistemy razrabotki/SRK, Moskva. -2021. -73 s. {[}in
Russian{]}

2. Abramkin N.I., Efimov V.I., Mansurov P.A. Jempiricheskie metodiki
ocenki sostojanija massiva gornyh porod// Izvestija UGGU. -2021. -T.
4(64). -S.109-115. DOI 10.21440/2307--2091-2021-4-109-115. {[}in
Russian{]}

3. Deere D.U., Deere D.W. The Rock Quality Designation (RQD) Index in
Practice/Rock Classification Systems for Engineering Purposes. ASTM STP
984. -1988. ISBN 0-8031-0988-1.

4. Bieniawski Z.T. Engineering Rock Mass Classifications: A Complete
Manual for Engineers and Geologists in Mining, Civil, and Petroleum
Engineering// New York: Wiley. -1989. ISBN:~9780471601722, 0471601721.

5. Grimstad E., Barton N. Updating of the Q-system for the Norwegian
Method of Tunnelling// Engineering Classification of Rock Masses for the
Design of Tunnel Support. Rock Mech. -1993. -P.189-239.

6. Hoek E., Carranza-Torres C., Corkum B., et al. Hoek--Brown Failure
Criterion - 2002 Edition // Proceedings of NARMS-TAC conference. -2002.
-P.267--273.

7. Laubscher D.H. A Geomechanics Classification System for the Rating of
Rock Mass in Mine Design// Journal of the South African Institute of
Mining and Metallurgy. -1990. -Vol.90(10). -P.257-273. DOI
10.10520/AJA0038223X\_1954.

8. Marinos P., Hoek E. Geological Strength Index: A Geologically Friendly
Tool for Rock Mass Strength Estimation// Proceedings of GeoEng 2000.
-2000. -P.1422-1446.

9. Hoek E., Diederichs M.S. Empirical Estimation of Rock Mass Modulus //
International Journal of Rock Mechanics and Mining Sciences. -2006.
-Vol.43(2). -P.203-215. DOI 10.1016/j.ijrmms.2005.06.005.

10. Polemis Júnior K., Chagas da Silva Filho F., Pinheiro Lima-Filho F.
Estimating the Rock Mass Deformation Modulus: A Comparative Study of
Empirical Methods Based on 48 Rock Mass Scenarios// REM - International
Engineering Journal. -2021. -Vol.74(1). -P.39-49. DOI
10.1590/0370-44672019740150.

11. Eremenko V.A., Ajnbinder I.I., Packevich P.G., Babkin E.A. Ocenka
sostojanija massiva gornyh porod na rudnikah ZF OAO "GMK
"Noril' skij nikel'" // Gornyj
informacionno-analiticheskij bjulleten'. -2017. -№1.
-C.5-17. ISSN 0236-1493. {[}in Russian{]}

12. Feng X.T., Gong B., Liang Z.,~\emph{et al.} tudy of the Dynamic
Failure Characteristics of Anisotropic Shales Under Impact Brazilian
Splitting // Rock Mechanics and Rock Engineering. -2024. -Vol.57. -P.
2213-2230. DOI 10.1007/s00603-023-03673-w.

13. Wen S., et al. Mechanical behavior of layered rocks under dynamic
loading // International Journal of Rock Mechanics and Mining Sciences.
-2023. -Vol.170:105533. DOI 10.1016/j.ijrmms.2023.105533.

14. Qiang Y., et al. Study on closing and cracking stress calculation
method of fractured rock //Frontiers in Earth Science. -2023. -Vol.
10:839304. DOI 10.3389/feart.2022.839304.

15. Huang L., et al. Mathematical problems in rock mechanics and
stress--strain evolution under thermal effects // Mathematics. -2023.
-Vol.11(1):67. DOI 10.3390/math11010067.

16. Mamaev Ju.A., Vlasov A.N., Mnushkin M.G., Jastrebov A.A. Izuchenie
naprjazhennogo sostojanija i mehanizma deformirovanija massivov gornyh
porod pri obrazovanii prirodno-tehnogennyh karstovyh provalov//
Geojekologija. Inzhenernaja geologija. Gidrogeologija. Geokriologija.
-2019. -T.1. -P.46-59. DOI 10.31857/S0869-78092019146-59. {[}in
Russian{]}

17. Potvin Y. Empirical Open Stope Design in Canada//The University of
British Columbia. - 1988. DOI 10.14288/1.0081130.

18. Brady B.H.G., Brown E.T. Rock Mechanics for Underground Mining.3rd
ed./ Dordrecht: Springer. -2006. DOI 10.1007/978-1-4020-2116-9.

19. Hutchinson D.J., Diederichs M.S. Cablebolting in Underground Mines/
BiTech Publishers. -1996. -406 p. ISBN 9780921095378.

20. Kaiser P.K., Cai M. Design of Rock Support System under Rockburst
Condition// Journal of Rock Mechanics and Geotechnical Engineering.
-2012. -4(3). -P.215-227. DOI 10.3724/SP.J.1235.2012.00215.

21. Abramkin N.I., Dorodnij A.V., Buharbaev I.U.Analiz integrirovannoj
tehnologii vysokoproizvoditel' noj otrabotki zapasov
vyemochnyh uchastkov ugol' nyh
shaht//Ugol'. -2019. -№1(1114). -S.40-45. DOI
10.18796/0041-5790-2019-1-40-45. {[}in Russian{]}

22. Kane W.F. Monitoring Slope Movement with Time Domain Reflectometry//
Engineering, Environmental Science, Geology. -2000.

23. Sarr D., Sall O., Ndiaye M., Lompo N. Comparative Analysis of the
Hard Hillsides Stability by Empirical Methods and Limit Equilibrium:
Case of Ultra Basic and Andesites of Mako and Marbles of Bandafassi
(Senegal)// Geomaterials. -2019. -Vol.9(3). -P.67-79. DOI
10.4236/gm.2019.93006.

24. Haimson B.C., Cornet F.H. ISRM Suggested Methods for Rock Stress
Estimation - Part 3: Hydraulic Fracturing (HF) and/or Hydraulic Testing
of Pre-existing Fractures (HTPF)// International Journal of Rock
Mechanics and Mining Sciences. -2003. -Vol.40(7-8). -P.1011-1020. DOI
10.1016/j.ijrmms.2003.08.002.

25. Kuranov A.D., Bagautdinov I.I. Experience of Geomechanical Research
and Calculations toward Ore Deposit Development in Complex Mining
Conditions// ~ISRM European Rock Mechanics Symposium - EUROCK 2018.
-2018.
\end{refs}

\begin{info}
\hspace{1em}\emph{{\bfseries Information about the authors}}

Shaposhnik S. N. - Scientific Supervisor, Limited Liability Partnership
``EXPERT PRO'', Ust-Kamenogorsk, Kazakhstan, e-mail:
shaposhniksergey@mail.ru;

Shokarev D. A. - Advisor to the Director for Strategic Affairs, Limited
Liability Partnership ``EXPERT PRO'', Ust-Kamenogorsk, Kazakhstan,
e-mail: D.Shokarev@expertpro.kz

Matvienko N.V. - Head of R\&D, Limited Liability Partnership ``EXPERT
PRO'', Ust-Kamenogorsk, Kazakhstan, e-mail:
n.matvienko@expertpro.kz;

Zikirova K. T. - Mineral Processing Engineer, Limited Liability
Partnership ``EXPERT PRO'', Ust-Kamenogorsk, Kazakhstan, e-mail:
k.zikirova@expertpro.kz;

Dauletzhanova Zh. T. - PhD Mining Engineering, Technology, Leading
Researcher of LLP "Institute of Coal Chemistry and Technology", K.
Kulazhanov Kazakh University of Technology and Business, Astana,
Kazakhstan, e-mail:
kaliyeva\_zhanna@mail.ru;

Konurin A. I. - Senior Researcher, Federal State Budgetary Institution
of Science, N.A. Chinakal Institute of Mining, Siberian Branch of the
Russian Academy of Sciences (SB RAS), Novosibirsk, Russia, e-mail:
Konurin@misd.ru.

\hspace{1em}\emph{{\bfseries Сведения об авторах}}

Шапошник С. Н. - научный руководитель, ТОО «EXPERT PRO», Астана,
Казахстан, e-mail: shaposhniksergey@mail.ru;

Шокарев Д. А. - советник директора по стратегическим вопросам, ТОО
«EXPERT PRO», Астана, Казахстан, e-mail:
D.Shokarev@expertpro.kz;

Матвиенко Н. В. - руководитель отдела НИОКР, ТОО «EXPERT PRO», Астана,
Казахстан, e-mail:
n.matvienko@expertpro.kz;

Зикирова К.Т. - инженер по обогащению полезных ископаемых, ТОО «EXPERT
PRO», Астана, Казахстан, e-mail:
k.zikirova@expertpro.kz;

Даулетжанова Ж.Т. - PhD, горное дело и технологии, ведущий научный
сотрудник ТОО «Институт химии угля и технологии», Казахский университет
технологии и бизнеса им/ К. Кулажанова, Астана, Казахстан, e-mail:
kaliyeva\_zhanna@mail.ru;

Конурин А. И. - старший научный сотрудник, ФГБУН «Институт горного дела
им. Н.А. Чинакала» СО РАН, Новосибирск, Россия, e-mail: Konurin@misd.ru.
\end{info}
